\documentclass[11pt,a4paper]{article}

\usepackage[utf8]{inputenc}
\usepackage[T1]{fontenc}
\usepackage[margin=2.5cm]{geometry}
\usepackage{amssymb}
\usepackage{parskip}
\usepackage{microtype}
\usepackage{hyperref}
\usepackage{xurl}
\usepackage{csquotes}
\usepackage{enumitem}
\usepackage{fancyvrb}

\hypersetup{
  colorlinks=true,
  linkcolor=black,
  urlcolor=blue,
  citecolor=blue,
  breaklinks=true
}

\emergencystretch=3em
\sloppy

\setlist[itemize]{label=$\blacktriangleright$}

\newcommand{\code}[1]{\texttt{#1}}
\newcommand{\pkg}[1]{\textsf{#1}}

\title{Response to JSS Editor's Preliminary Review\\
\large \textit{BayesianDEB}: A Bayesian Framework for Dynamic Energy
Budget Modelling in R}
\author{Branimir K.\ Hackenberger\\ \small on behalf of the authors}
\date{2026-05-11}

\begin{document}
\maketitle

We thank the editor for the close reading of \emph{BayesianDEB: A
Bayesian Framework for Dynamic Energy Budget Modelling in R}. The
comments group naturally into three areas: the R class system, the
replication material, and the package vignettes. We have worked through
all of them. Each comment is quoted below, word for word or as a faithful
summary, followed by our reply and the commits and files it touches in
version \textbf{0.2.0} of the package.

The work lives on the \code{jss-revision} branch of
\url{https://github.com/sciom/BayesianDEB} and is tagged \code{v0.2.0}.
The final package is \code{BayesianDEB\_0.2.0.tar.gz} and the replication
archive is \code{BayesianDEB\_replication.zip}.

\section{Class system}

\subsection{\code{bdeb\_data} and \code{bdeb\_model} are missing
\code{summary()} and \code{plot()}}

\begin{displayquote}
``\verb|methods(class = "bdeb_data")| shows only \code{print}.
Journal of Statistical Software requires that at least \code{summary}
and \code{plot} are also implemented for every class in the package.
The same remark holds for \code{bdeb\_model}.''
\end{displayquote}

\textbf{Response.} Added. \code{bdeb\_data}, \code{bdeb\_model} and
\code{bdeb\_prior} now carry a complete
\code{print()}~/ \code{summary()}~/ \code{plot()} trio. Each
\code{summary()} method returns a tidy \code{summary.bdeb\_data}~/
\code{summary.bdeb\_model}~/ \code{summary.bdeb\_prior} object with its
own \code{print()} method, and each \code{plot()} method gives an
overview of the data, the model structure, or the prior densities.

\textbf{Reference.} Tasks B5/B6/B7 of the revision plan; commits in
\code{R/data\_prep.R}, \code{R/model\_spec.R}, \code{R/priors.R}.
Tests added in \code{tests/testthat/test-api-contracts.R}.

\subsection{\code{bdeb\_diagnose()} returns a plain \code{list}}

\begin{displayquote}
``\code{bdeb\_diagnose} prints a very long information and
returns an object of class \code{list}. This should be improved. It
would be better that \code{bdeb\_diagnose} returns a specific class
object with methods to inspect it and that its \code{print} is the
output of the \code{print} method for this class.''
\end{displayquote}

\textbf{Response.} \code{bdeb\_diagnose()} now returns a
\code{bdeb\_diagnostics} S3 object with \code{print()},
\code{summary()} and \code{plot()} methods. It no longer prints to the
console as a side effect: the output comes from \verb|print(d)|, or from
autoprinting at the top level. The \code{plot()} method takes
\verb|type = "rhat"| and \verb|type = "ess"|, and the old list-style
access (\verb|d$n_divergent|, \verb|d$summary|) still works for backward
compatibility.

\textbf{Reference.} Task B5; \code{R/diagnostics.R}.

\subsection{\code{predict(fit)} returns an unstructured long list}

\begin{displayquote}
``\verb|predict(fit)| is a long list with multiple entries. We
would expect a specific class for this object as well, to improve
readability of the output for users and allow the implementation of
methods that would allow inspecting this object.''
\end{displayquote}

\textbf{Response.} \code{bdeb\_predict()}, the function behind
\code{predict.bdeb\_fit}, now returns a \code{bdeb\_prediction} S3 object
with \code{print()}, \code{summary()} and \code{plot()} methods.
\code{print()} reports the model type, the prediction horizon and the
posterior summary widths; \code{summary()} hands back a tidy data frame
with \code{time}~/ \code{lower}~/ \code{median}~/ \code{upper} columns.

\textbf{Reference.} Task B6; \code{R/fit.R}, \code{R/plot.R}.

\subsection{\code{bdeb\_fit} has too few methods}

\begin{displayquote}
``\verb|methods(class = "bdeb_fit")| returns
\verb|[coef plot predict print summary]| which seems to be a bit
limited. For model classes, we would expect more methods allowing to
inspect the output and assess the quality of the results (see, e.g.,
\verb|methods(class = "lm")|), including \code{confint},
\code{residuals}, \ldots''
\end{displayquote}

\textbf{Response.} We added the \code{lm}-style methods that make sense
for a Bayesian DEB fit:

\begin{itemize}[nosep]
\item \code{confint()}: posterior credible intervals.
\item \code{fitted()}: posterior median/mean of $\hat{L}_i$.
\item \code{residuals()}: observed minus fitted.
\item \code{nobs()}: observation count.
\item \code{vcov()}: posterior covariance of model parameters.
\item \code{logLik()}: log-pointwise predictive density
  (\code{lppd}).
\end{itemize}

\code{fitted()}, \code{residuals()} and \code{logLik()} apply only to the
\verb|"individual"| and \verb|"growth_repro"| model types; on
\verb|"hierarchical"| and \verb|"debtox"| fits, where a
single-observation residual has no clear meaning, they stop with an
explanatory error.

\textbf{Reference.} Tasks B10/B11; \code{R/fit.R}.

\subsection{\code{bdeb\_derived()} should be a method on
\code{bdeb\_fit}}

\begin{displayquote}
``\code{bdeb\_derived} starts with
\verb|if (!inherits(fit, "bdeb_fit"))| which shows that it could have
been implemented as a method rather than a simple function to clarify
the class structure for the user.''
\end{displayquote}

\textbf{Response.} \code{bdeb\_derived()} is now an S3 generic, with
\code{bdeb\_derived.bdeb\_fit} as its method. The call users write is
unchanged, but dispatching this way makes the class structure visible and
leaves room for methods on prior-only and simulated objects later.

\textbf{Reference.} Task B7; \code{R/fit.R}.

\subsection{\code{summary(fit)} and \code{bdeb\_summary(fit)} are
redundant}

\begin{displayquote}
``\verb|summary(fit)| and \verb|bdeb_summary(fit)| seem to be
the same thing, which questions the reason why the second has been
implemented.''
\end{displayquote}

\textbf{Response.} \code{summary.bdeb\_fit()} is now the main way to
summarise the posterior: it takes \code{pars} and \code{prob} arguments
and returns a \code{posterior::draws\_summary} data frame, in the spirit
of \code{summary.lm()}. \code{bdeb\_summary()} is deprecated. The wrapper
still runs, but it raises a \pkg{lifecycle} deprecation warning on first
use and will go in a future release.

\textbf{Reference.} Task B9; \code{R/fit.R}, \code{R/diagnostics.R}.

\subsection{Input validation is incomplete}

\begin{displayquote}
``Functions do not always check their inputs. All inputs must
be checked to assess that their value is what is expected and returns
comprehensive messages to users when this is not the case.''
\end{displayquote}

\textbf{Response.} Every exported function now checks its inputs through
one shared set of internal assertions (\code{assert\_positive},
\code{assert\_finite\_scalar}, \code{assert\_probability}, and so on).
Error messages go through \code{cli::cli\_abort()} with named bullets and
suggested fixes. Input validation is now covered by more than 180
contracts; \code{tests/testthat/test-input-validation.R} (63
\code{expect\_error} checks) and \code{test-api-contracts.R} (38 checks)
between them exercise the entry points users hit most.

\textbf{Reference.} Task B8; \code{R/utils.R} for the assertions,
individual R files for their use.

\section{Replication material}

\subsection{Replication material organisation and README}

\begin{displayquote}
``A single zip folder should be submitted which contains all
scripts and results. If more than one script needs to be provided, a
README should be added which explains which part of the replication
material reproduces which part of the manuscript.''
\end{displayquote}

\textbf{Response.} The new \code{replication/} directory ships as a
single zip (\code{BayesianDEB\_replication.zip}), laid out like this:

\begin{Verbatim}[fontsize=\small]
replication/
  README.md                # one-page guide: what each script does
  00_setup.R               # libraries, palette, BDEB_MODE
  01_illustrations.R       # Section 5 of the manuscript
  02_validation.R          # Section 6
  03_case_studies.R        # Section 7
  comparison_debinfer.R    # Section 9 auxiliary comparison
  data/                    # bundled inputs (curves.txt, ...)
  outputs/                 # generated figures and tables
  sbc/                     # long-running SBC scripts
\end{Verbatim}

\code{README.md} maps each script to the manuscript sections it produces,
gives the run time in both \code{lite} and \code{full} modes, and lists
the required and optional packages.

\textbf{Reference.} Phase 3 (tasks A1--A7); \code{replication/README.md}.

\subsection{Manuscript code missing from replication}

\begin{displayquote}
``All code shown in the manuscript should also be included in
the replication scripts. E.g., page 7 contains
\verb|prior_species("Eisenia_fetida")|. However, we were unable to
find this code in the replication scripts provided.''
\end{displayquote}

\textbf{Response.} Every code listing in the manuscript is now run in the
replication scripts. \verb|prior_species("Eisenia_fetida")| is called in
\code{replication/01\_illustrations.R} just before the matching
\code{bdeb\_fit()} call from Section~5.1, and the \code{obs\_student\_t()}
one-liner from p.~18 is now constructed in Section~5.1 as well, so the
audit is complete, even though nothing is fitted with it there.

\textbf{Reference.} \code{replication/01\_illustrations.R} lines for
\code{prior\_species()} and \code{obs\_student\_t()}.

\subsection{Differences between replication output and manuscript}

\begin{displayquote}
``Running replication material shows that there are slight
differences between the outputs produced by the replication material
and what is shown in the manuscript. Please fix this or explain why
identical results cannot be obtained by specifying the specific
environment that allowed to obtain those results in the manuscript.''
\end{displayquote}

\textbf{Response.} Those differences came from two things in the earlier
scripts: two figures used an unfixed seed, and the exact Stan, cmdstanr
and R versions behind the manuscript were not recorded. We have:

\begin{itemize}[nosep]
\item fixed a global seed (\verb|set.seed(20260418)|) in
  \code{00\_setup.R} and pass it to every \code{bdeb\_fit()} and
  \code{plot\_*()} call through explicit \verb|seed =| arguments;
\item recorded the rendering environment in
  \code{replication/outputs/sessionInfo.txt}, captured with
  \code{BayesianDEB::bdeb\_session\_info()} and
  \code{utils::sessionInfo()};
\item with this environment, \verb|BDEB_MODE=full Rscript ...| reproduces
  every manuscript figure and table bit for bit.
\end{itemize}

A small residual difference ($\sim 10^{-4}$ in the posterior summaries)
still shows up across platforms, because Stan's adaptive HMC depends on
the order of floating-point reductions in \code{reduce\_sum}; we note
this in \code{replication/README.md}.

\textbf{Reference.} Tasks A3/A4; \code{replication/00\_setup.R},
\code{replication/outputs/sessionInfo.txt}.

\subsection{CVode mxstep informational messages}

\begin{displayquote}
``The replication material repeatedly produces these warnings:
\code{CVode(...\ CV\_NORMAL)\ failed\ with\ error\ flag\ -1:\ The
solver took mxstep internal steps but could not reach tout.} That
seems to be important to fix or at least provide an explanation and
guide users in how to interpret these warnings.''
\end{displayquote}

\textbf{Response.} Stan's stiff BDF ODE solver prints these during HMC
warm-up, when the sampler wanders into parameter combinations that do not
fit the data. They are \emph{informational} rather than errors, and Stan
documents them as such: the integrator hands control back, the leapfrog
step is rejected by the Metropolis criterion, and the chain carries on.

We made two changes:

\begin{enumerate}[nosep]
\item Increased \code{ode\_bdf\_tol} \code{max\_num\_steps} from
  \code{1e4} to \code{1e5} in all four Stan models
  (\code{inst/stan/bdeb\_*.stan}). This removes roughly two-thirds of the
  messages on the default 300-iteration warm-up.
\item Added a paragraph to \code{replication/README.md} and to the
  \emph{Getting started} vignette (``ODE solver messages'') that quotes
  the Stan reference manual and explains when the messages point to a
  real problem (they recur past warm-up, alongside divergent
  transitions) and when they are just harmless warm-up exploration
  (sporadic, warm-up only).
\end{enumerate}

\textbf{Reference.} Task A4 / C1; \code{inst/stan/bdeb\_*.stan},
\code{replication/README.md}, \code{vignettes/getting\_started.Rmd}.

\subsection{External \code{curves.txt} dependency}

\begin{displayquote}
``The replication material contains
\verb|curves <- read.delim(|\allowbreak\verb|"https://raw.githubusercontent.com/|\allowbreak
\verb|JeromeMathieuEcology/|\allowbreak\verb|EGrowth/master/curves.txt")|.
This file should also be directly included in the replication material
to reduce the dependency on availability from external unreliable
sources.''
\end{displayquote}

\textbf{Response.} \code{curves.txt} now travels with the material at
\code{replication/data/curves.txt} (SHA-256 \code{19e131b8\ldots},
recorded in \code{replication/data/CHECKSUMS.txt}). The scripts read the
local copy, and the URL-based form is gone. Its provenance and licence
(CC-BY-4.0, Mathieu 2024) are noted in \code{replication/README.md}.

\textbf{Reference.} Task A5; \code{replication/data/curves.txt},
\code{replication/data/CHECKSUMS.txt}.

\subsection{\code{eisenia\_real\_data.R} fails with \code{T~=~298.15}}

\begin{displayquote}
``Running \code{eisenia\_real\_data.R} fails with
\verb|Error in bdeb_model():|\allowbreak{} \verb|Temperature list missing: T_obs.|''
\end{displayquote}

\textbf{Response.} In version 0.1.4 we renamed the temperature-list field
\code{T} to \code{T\_obs}, both to match the Stan-side name and to stop it
shadowing R's built-in \code{T~=~TRUE}. The replication scripts now
follow suit: every \verb|temperature = list(...)| call uses
\verb|T_obs = ...|. We weighed a backward-compatibility shim but left it
out to keep the v0.2.0 API clean.

\textbf{Reference.} Task A6; \code{replication/03\_case\_studies.R}.

\subsection{Computing resource budget}

\begin{displayquote}
``We would expect that the replication material runs within a
reasonable amount of time using standard computing resources.
Excessive use of computing resources should be avoided when providing
examples to demonstrate the use of the package.''
\end{displayquote}

\textbf{Response.} Each replication script runs in one of two modes,
chosen with the \code{BDEB\_MODE} environment variable:

\begin{itemize}[nosep]
\item \textbf{\code{lite}} (default): \verb|chains = 2|,
  \verb|iter_warmup = 300|, \verb|iter_sampling = 300|. Reproduces the
  qualitative shape of every figure and table. Total wall-clock on a
  4-core laptop: \textbf{< 30 minutes} for \code{01} + \code{02} +
  \code{03}.
\item \textbf{\code{full}} (manuscript-equivalent): \verb|chains = 4|,
  \verb|iter_warmup = 1000|, \verb|iter_sampling = 1000|. Reproduces
  bit-identical posteriors and diagnostics. Total wall-clock on a 4-core
  laptop: \textbf{$\approx 3$ hours}.
\end{itemize}

\code{README.md} documents the choice. The SBC calibration scripts
(\code{replication/sbc/}) stay in their own sub-folder: they are
long-running by nature ($\sim 10$~h) and are not needed to reproduce the
manuscript figures, so their results ship as cached \code{.rds} files.

\textbf{Reference.} Task A2; \code{replication/00\_setup.R} (mode
dispatch), \code{replication/README.md}.

\section{Vignettes and examples}

\subsection{\code{bdeb\_model} has no examples}

\begin{displayquote}
``\verb|example("bdeb_model")| warns
\verb|'bdeb_model' has a help file|\allowbreak{} \verb|but no examples|. Since
\code{bdeb\_model} seems to be one of the main functions of the
package, it is important that its help page is accompanied with an
example.''
\end{displayquote}

\textbf{Response.} We added runnable examples to every main function
whose example does not need CmdStan: \code{bdeb\_data},
\code{bdeb\_model}, \code{bdeb\_tox}, \code{bdeb\_prior\_predictive},
\code{prior\_species}, and all the \code{prior\_*} constructors. Each runs
on bundled datasets in well under a second under R~CMD~check defaults,
which we confirmed with \verb|R CMD check --as-cran --run-donttest|.

\textbf{Reference.} Task C2; \code{R/*.R}
(Roxygen \code{@examples} blocks).

\subsection{\code{bdeb\_fit} examples in \code{\textbackslash dontrun\{\}}}

\begin{displayquote}
``\verb|example(bdeb_fit)| shows that all code is in
\code{dontrun}. It would seem preferable to conditionally run the
code after checking that the external CmdStan toolchain is available,
similar to what is suggested for suggested packages in Writing R
Extensions.''
\end{displayquote}

\textbf{Response.} \code{bdeb\_fit()} keeps its
\code{\textbackslash dontrun\{\}} example for two reasons: compiling a
Stan model and then running MCMC goes well past the 5-second-per-example
limit in R~CMD~check, and the CmdStan toolchain has to be installed and
tested deliberately rather than skipped on failure. The Roxygen block now
says this in a comment that shows up in the rendered Rd, so a reader can
see why the example is wrapped.

For coverage we lean on two full vignettes (\emph{Getting started};
\emph{Case study: Eisenia / Folsomia}) and the replication scripts, all
of which guard their Stan chunks with
\verb|requireNamespace("cmdstanr", quietly = TRUE)|.

\textbf{Reference.} Task C2 (Roxygen comment); Task C1 (vignette
conditionality, see~3.3).

\subsection{Vignette code commented out}

\begin{displayquote}
``All code in vignettes seems to be commented out, see for
example
\url{https://cran.r-project.org/web/packages/BayesianDEB/vignettes/case_study_eisenia_folsomia.R}.
It would seem preferable to check for availability of \pkg{cmdstanr}
and then conditionally allow to run the code.''
\end{displayquote}

\textbf{Response.} Both vignettes now gate every Stan-touching chunk on
\verb|requireNamespace("cmdstanr", quietly = TRUE)|, following the
\emph{Writing R Extensions} advice for suggested packages: where
\pkg{cmdstanr} is installed \emph{and} \code{cmdstanr::cmdstan\_version()}
returns a non-\code{NULL} version string, the chunk runs and shows the
posterior; otherwise it is skipped and a message explains how to install
CmdStan.

To keep the rendered HTML small, each fit drops to \code{lite} mode
(\verb|chains = 2, iter = 300+300|) when the build happens on CRAN; on a
full installation the same code still draws every figure from the
article.

\textbf{Reference.} Task C1; \code{vignettes/getting\_started.Rmd},
\code{vignettes/case\_study\_eisenia\_folsomia.Rmd}.

\subsection{\code{plot(fit, type = "pairs")} returns NULL}

\begin{displayquote}
``Running the code it would seem that no plot is created with:
\verb|plot(fit, type = "pairs", pars = c("p_Am", "p_M", "kappa"))|
\code{NULL}''
\end{displayquote}

\textbf{Response.} The old \verb|plot.bdeb_fit(type = "pairs")| called
\code{bayesplot::mcmc\_pairs()}, which returns a \code{bayesplot\_grid}
(a \code{gtable}), and then tried \verb|result + ggplot2::labs()|, which
is undefined for \code{gtable} objects and quietly returned \code{NULL}.

The method now returns the \code{bayesplot\_grid} directly, and the
requested parameters reach \code{bayesplot::mcmc\_pairs()} through
\code{posterior::subset\_draws()}, so the \code{pars} argument is honoured
exactly. Tests in \code{tests/testthat/test-plot.R} check that the return
value is non-NULL and that the grid has the right dimensions.

\textbf{Reference.} Task C1; \code{R/plot.R}.

\section{Additional improvements (not requested, but bundled)}

Since the class-system rework reached into most of the package, we used
the occasion to:

\begin{itemize}[nosep]
\item Add LOO and WAIC predictive-density comparison
  (\code{bdeb\_loo()}).
\item Add the \code{obs\_student\_t()} observation family for robust
  growth fits.
\item Add Arrhenius temperature correction to all four Stan models
  (previously only \verb|"individual"|).
\item Add within-chain parallelism via \code{reduce\_sum} to the
  \verb|"hierarchical"| and \verb|"debtox"| models, with the
  \code{threads\_per\_chain} argument exposed in \code{bdeb\_fit()}.
\item Grow the test suite from 58 \code{test\_that} blocks in version
  0.1.0 to 395 blocks across 18 files (API contracts, scientific
  consistency, snapshot regression, deep validation, end-to-end
  integration, input validation, class methods), with over 180
  input-validation contracts as detailed in~1.7.
\end{itemize}

All of these are recorded in \code{NEWS.md} under the 0.2.0 entry.

\section{Submission summary}

The revised submission consists of:

\begin{itemize}[nosep]
\item \code{BayesianDEB\_0.2.0.tar.gz}: the package, passes
  \verb|R CMD check --as-cran| with \textbf{0 errors / 0 warnings /
  1 note} (the single NOTE is the informational ``unable to verify
  current time'', produced when the host has no network access to
  the time server; it is unrelated to package contents).
\item \code{BayesianDEB\_replication.zip}: replication material,
  with \code{README.md}, bundled \code{curves.txt}, and \code{lite}
  / \code{full} modes.
\item The present document (\code{response\_to\_editor\_round1.tex}).
\end{itemize}

We thank the editor and reviewers for the thorough, constructive
feedback, which has improved the package considerably.

\vspace{1em}
\noindent Branimir K.\ Hackenberger\\
\noindent on behalf of the authors\\
\noindent 2026-05-11

\end{document}
